\section{Introduction}\label{sec:intro}

\begin{figure*}[t]
  \centering
  \includegraphics[width=\textwidth]{20_snap_hyre_pipeline_art.png}
  \caption{\hyre{}: the first model call writes a private draft answer and a
  separate legal search passage; only the search passage is used to retrieve
  evidence. The second model call sees only the original question and the
  top-$k$ retrieved passages, and produces the final answer. The draft is
  never passed to the second call.}
  \label{fig:pipeline_art}
\end{figure*}

Retrieval-Augmented Generation (RAG) has emerged as a powerful paradigm for enhancing Large Language Model (LLM) generation by incorporating external evidence retrieved during inference. RAG is particularly well-suited for legal question answering, as responses frequently depend on statutes, regulations, court holdings, and doctrinal rules. This domain-specific information typically falls outside the core knowledge base of LLMs, particularly reasoning models.
A standard RAG pipeline typically generates an initial query to retrieve relevant external evidence, which is then concatenated into the LLM's context \citep{lewis2020rag}. However, formulating an effective initial query is often the most challenging step. For example, a bar-exam fact pattern might emphasize concrete narrative details such as a driver, crutches, and a banana peel, whereas the necessary supporting passage articulates an abstract, well-established rule regarding proximate cause. Therefore, the standard RAG is prone to generating inappropriate or superficial queries for complex legal facts, leading to retrieval failures.

A common solution is to let the model first \emph{write} the kind of passage it expects to find and search with this text instead of the raw question. HyDE (Hypothetical Document Embeddings) prompts the model to write a passage to answer the question, and uses its embedding to retrieve actual passages~\citep{gao2023hyde}. Similarly, query expansion rewrites a pseudo-document as the query for search~\citep{wang2023query2doc}. While improving retrieval accuracy in general domains, they present significant challenges in legal contexts, where colloquial fact patterns should be mapped to formal statutory language. In specific legal cases, the reformulation step frequently breaks down. Consequently, they struggle to find correct statutory expressions, often abstracting away critical anchors, applying the wrong legal frame, or retrieving plausible but irrelevant authority.

To address that issue, we introduce \hyre{} (Snap-answer Conditioned Pseudo-document Embedding), a decoupled, two-stage query generation method designed to simulate expert legal reasoning. When facing a complex fact pattern, a legal practitioner typically forms a quick initial intuition, uses that intuition to locate controlling authority, and subsequently evaluates the actual evidence on its own merits rather than anchoring on their first guess. \hyre{} structurally encodes this cognitive loop. In the first phase, the model is prompted to generate a ``snap answer'', a rapid, unverified draft response reflecting its immediate legal intuition based purely on internal knowledge. Conditioned directly on this initial snap answer, the model then generates a pseudo-document written in the formal, abstract style of the supportive legal authority that would resolve the case. Crucially, during the retrieval phase, the snap answer is discarded to prevent confirmation bias. The pseudo-document is embedded to retrieve actual legal evidence from the corpus. In the final answer phase, a separate model call generates the final answer conditioned on the retrieved external evidence and the original question, ensuring the final reasoning is grounded in real authority and remains uncontaminated from the noisy initial guess.

We evaluate two benchmarks selected to represent opposite ends of the question-corpus vocabulary gap~\cite{Zheng_2025}. The first benchmark, BarExamQA, exemplifies the \textit{weak-query regime}: its questions consist of narrative fact patterns detailing accidents, transactions, and procedures in everyday prose, whereas the target corpus relies on formal doctrinal vocabulary (e.g., ``proximate cause,'' ``hearsay exception,'' or ``adverse possession''). Consequently, the raw question and the relevant legal rule share extremely low lexical similarity. In contrast, HousingQA represents the \textit{strong-query regime}. These statutory-entailment questions inquire directly about specific legal conditions (e.g., ``Can a landlord evict a tenant for committing waste?''). Because the queries explicitly utilize formal terminology, their vocabulary naturally overlaps with the target housing statutes required to answer them.

On the BarExamQA benchmark, where questions show a large gap with the corpus, \hyre{} yields substantial improvements. Across the 8B, 26B, and 70B model sizes, it lifts the gold passage Hit@5 rate from a baseline of 1.4\% to 9.5\%-12.1\%, consistently improving final answer accuracy at every model scale.

In summary, we formalize \hyre{} as a decoupled, two-stage RAG framework, which is highly tailored for Legal RAG challenges. Crucially, when user questions describe scenarios using narrative prose that lacks lexical overlap with formal legal texts, \hyre{} successfully bridges this gap to drive major improvements in gold passage retrieval, subsequently boosting final answer accuracy.
